% Chapter 6: Sard's Theorem
% Lee, *Introduction to Smooth Manifolds* (2nd ed., GTM 218).
% Generated from the section extraction; nodes carry only Lean that is
% verified sorry-free. See ../../UPSTREAM_LEAN_AUDIT.md.


\section{Sets of Measure Zero}

\begin{definition}
\label{def:6-38-extra-1}
\lean{NullMeasurableSet.of_null}
\mathlibok
Recall what it means for a set $A \subseteq  {\mathbb{R}}^{n}$ to have measure zero (see Appendix C): for any $\delta  > 0, A$ can be covered by a countable collection of open rectangles, the sum of whose volumes is less than $\delta$ (Fig. 6.1).
\end{definition}

\begin{exercise}
\label{exer:6-1}
\lean{volume_eq_zero_iff_forall_pos_exists_open_cube_cover}
\leanok
\dcref{ch6:6.1}
Show that open rectangles can be replaced by open balls or open cubes in the definition of subsets of measure zero. [Hint: for cubes, first show that every open rectangle $R \subseteq  {\mathbb{R}}^{n}$ can be covered by finitely many open cubes, the sum of whose volumes is no more than ${2}^{n}$ times the volume of $R$ . For balls, Fig. 6.2 suggests the main idea.]
\end{exercise}

\begin{lemma}
\label{lem:6-2}
\lean{volume_eq_zero_of_isCompact_of_forall_first_coordinate_slice_volume_eq_zero}
\leanok
\dcref{ch6:6.2}
Suppose $A \subseteq  {\mathbb{R}}^{n}$ is a compact subset whose intersection with $\{ c\}  \times  {\mathbb{R}}^{n - 1}$ has $\left( {n - 1}\right)$ -dimensional measure zero for every $c \in  \mathbb{R}$ . Then $A$ has n-dimensional measure zero.
\end{lemma}

\begin{proof}
\leanok
Choose an interval $\left\lbrack  {a, b}\right\rbrack   \subseteq  \mathbb{R}$ such that $A \subseteq  \left\lbrack  {a, b}\right\rbrack   \times  {\mathbb{R}}^{n - 1}$ . For each $c \in \; \left\lbrack  {a, b}\right\rbrack$ , let ${A}_{c} \subseteq  {\mathbb{R}}^{n - 1}$ denote the compact subset $\left\{  {x \in  {\mathbb{R}}^{n - 1} : \left( {c, x}\right)  \in  A}\right\}$ .

Let $\delta  > 0$ be given. Our hypothesis implies that for each $c \in  \left\lbrack  {a, b}\right\rbrack$ , the set ${A}_{c}$ is covered by finitely many $\left( {n - 1}\right)$ -dimensional open cubes ${C}_{1},\ldots ,{C}_{k}$ with total volume less than $\delta$ . Let ${U}_{c}$ be the open subset ${C}_{1} \cup  \cdots  \cup  {C}_{k} \subseteq  {\mathbb{R}}^{n - 1}$ . Because $A$ is compact, there must be an open interval ${J}_{c}$ containing $c$ such that the intersection of $A$ with ${J}_{c} \times  {\mathbb{R}}^{n - 1}$ is contained in ${J}_{c} \times  {U}_{c}$ , for otherwise there would be a sequence of points $\left( {{c}_{i},{x}_{i}}\right)  \in  A$ such that ${c}_{i} \rightarrow  c$ and ${x}_{i} \notin  {U}_{c}$ ; but then passing to a convergent subsequence we obtain ${x}_{i} \rightarrow  x \in  {A}_{c} \smallsetminus  {U}_{c}$ , which contradicts the fact that ${A}_{c} \subseteq  {U}_{c}$ .

The intervals $\left\{  {{J}_{c} : c \in  \left\lbrack  {a, b}\right\rbrack  }\right\}$ form an open cover of $\left\lbrack  {a, b}\right\rbrack$ , so there are finitely many numbers ${c}_{1} < \cdots  < {c}_{m}$ such that the intervals ${J}_{{c}_{1}},\ldots ,{J}_{{c}_{m}}$ cover $\left\lbrack  {a, b}\right\rbrack$ . By shrinking the intervals ${J}_{{c}_{i}}$ where they overlap if necessary, we can arrange that the combined lengths of ${J}_{{c}_{1}},\ldots ,{J}_{{c}_{m}}$ add up to no more than $2\left| {b - a}\right|$ . It follows that $A$ is contained in $\left( {{J}_{{c}_{1}} \times  {U}_{{c}_{1}}}\right)  \cup  \cdots  \cup  \left( {{J}_{{c}_{m}} \times  {U}_{{c}_{m}}}\right)$ , which is a union of finitely many open rectangles with combined volume less than ${2\delta }\left| {b - a}\right|$ . Since this can be made as small as desired, it follows that $A$ has $n$ -dimensional measure zero.
\end{proof}

\begin{proposition}
\label{prop:6-3}
\lean{volume_halfSpace_graph_eq_zero_of_measurableSet_of_continuous}
\leanok
\uses{lem:6-2}
\dcref{ch6:6.3}
Suppose $A$ is an open or closed subset of ${\mathbb{R}}^{n - 1}$ or ${\mathbb{H}}^{n - 1}$ , and $f : A \rightarrow  \mathbb{R}$ is a continuous function. Then the graph of $f$ has measure zero in ${\mathbb{R}}^{n}$ .
\end{proposition}

\begin{proof}
\leanok
First assume $A$ is compact. We prove the theorem in this case by induction on $n$ . When $n = 1$ , it is trivial because the graph of $f$ is at most a single point. To prove the inductive step, we appeal to Lemma 6.2. For each $c \in  \mathbb{R}$ , the intersection of the graph of $f$ with $\{ c\}  \times  {\mathbb{R}}^{n - 1}$ is just the graph of the restriction of $f$ to $\left\{  {x \in  A : {x}^{1} = c}\right\}$ , which is in turn the graph of a continuous function of $n - 2$ variables. It follows by induction that each such graph has $\left( {n - 1}\right)$ -dimensional measure zero, and thus by Lemma 6.2, the graph of $f$ itself has $n$ -dimensional measure zero.

If $A$ is noncompact, it is a countable union of compact subsets by Proposition A.60, so the graph of $f$ is a countable union of sets of measure zero.
\end{proof}

\begin{corollary}
\label{cor:6-4}
\lean{volume_affineSubspace_eq_zero_of_ne_top}
\leanok
\uses{prop:6-3}
\dcref{ch6:6.4}
Every proper affine subspace of ${\mathbb{R}}^{n}$ has measure zero in ${\mathbb{R}}^{n}$ .
\end{corollary}

\begin{proof}
\leanok
Let $S \subseteq  {\mathbb{R}}^{n}$ be a proper affine subspace. Suppose first that $\dim S = n - 1$ . Then there is at least one coordinate axis, say the ${x}^{i}$ -axis, that is not parallel to $S$ , and in that case $S$ is the graph of an affine function of the form ${x}^{i} = F\left( {{x}^{1},\ldots ,{x}^{i - 1},{x}^{i + 1},\ldots ,{x}^{n}}\right)$ , so it has measure zero by Proposition 6.3. If $\dim S < n - 1$ , then $S$ is contained in some affine subspace of dimension $n - 1$ , so it follows from Proposition C.18(b) that $S$ has measure zero.
\end{proof}

\begin{definition}
\label{def:6-38-extra-2}
\lean{has_measure_zero_in_manifold}
\leanok
If $M$ is a smooth $n$ -manifold with or without boundary, we say that a subset $A \subseteq  M$ has measure zero in $\mathbf{M}$ if for every smooth chart $\left( {U,\varphi }\right)$ for $M$ , the subset $\varphi \left( {A \cap  U}\right)  \subseteq  {\mathbb{R}}^{n}$ has $n$ -dimensional measure zero. The following lemma shows that we need only check this condition for a single collection of smooth charts whose domains cover $A$ .
\end{definition}

\begin{lemma}
\label{lem:6-6}
\lean{has_measure_zero_in_manifold_of_chart_cover}
\leanok
\dcref{ch6:6.6}
Let $M$ be a smooth $n$ -manifold with or without boundary and $A \subseteq  M$ . Suppose that for some collection $\left\{  \left( {{U}_{\alpha },{\varphi }_{\alpha }}\right) \right\}$ of smooth charts whose domains cover $A,{\varphi }_{\alpha }\left( {A \cap  {U}_{\alpha }}\right)$ has measure zero in ${\mathbb{R}}^{n}$ for each $\alpha$ . Then $A$ has measure zero in $M$ .
\end{lemma}

\begin{proof}
\leanok
Let $\left( {V,\psi }\right)$ be an arbitrary smooth chart. We need to show that $\psi \left( {A \cap  V}\right)$ has measure zero. Some countable collection of the ${U}_{\alpha }$ ’s covers $A \cap  V$ . For each such ${U}_{\alpha }$ , we have

\[
\psi \left( {A \cap  V \cap  {U}_{\alpha }}\right)  = \left( {\psi  \circ  {\varphi }_{\alpha }^{-1}}\right)  \circ  {\varphi }_{\alpha }\left( {A \cap  V \cap  {U}_{\alpha }}\right) .
\]

(See Fig. 6.4.)Now, ${\varphi }_{\alpha }\left( {A \cap  V \cap  {U}_{\alpha }}\right)$ is a subset of ${\varphi }_{\alpha }\left( {A \cap  {U}_{\alpha }}\right)$ , which has measure zero in ${\mathbb{R}}^{n}$ by hypothesis. By Proposition 6.5 applied to $\psi  \circ  {\varphi }_{\alpha }^{-1}$ , therefore, $\psi \left( {A \cap  V \cap  {U}_{\alpha }}\right)$ has measure zero. Since $\psi \left( {A \cap  V}\right)$ is the union of countably many such sets, it too has measure zero.
\end{proof}

\begin{exercise}
\label{exer:6-7}
\lean{has_measure_zero_in_manifold_iUnion}
\leanok
\dcref{ch6:6.7}
Let $M$ be a smooth manifold with or without boundary. Show that a countable union of sets of measure zero in $M$ has measure zero.
\end{exercise}

\begin{proposition}
\label{prop:6-8}
\lean{has_measure_zero_in_manifold}
\leanok
\dcref{ch6:6.8}
Suppose $M$ is a smooth manifold with or without boundary and $A \subseteq  M$ has measure zero in $M$ . Then $M \smallsetminus  A$ is dense in $M$ .
\end{proposition}

\begin{proof}
\leanok
If $M \smallsetminus  A$ is not dense, then $A$ contains a nonempty open subset of $M$ , which implies that there is a smooth chart $\left( {V,\psi }\right)$ such that $\psi \left( {A \cap  V}\right)$ contains a nonempty open subset of ${\mathbb{R}}^{n}$ (where $n = \dim M$ ). Because $\psi \left( {A \cap  V}\right)$ has measure zero in ${\mathbb{R}}^{n}$ , this contradicts Corollary C.25.
\end{proof}

\begin{theorem}
\label{thm:6-9}
\lean{has_measure_zero_in_manifold}
\leanok
\dcref{ch6:6.9}
Suppose $M$ and $N$ are smooth $n$ -manifolds with or without boundary, $F : M \rightarrow  N$ is a smooth map, and $A \subseteq  M$ is a subset of measure zero. Then $F\left( A\right)$ has measure zero in $N$ .
\end{theorem}

\begin{proof}
\leanok
Let $\left\{  \left( {{U}_{i},{\varphi }_{i}}\right) \right\}$ be a countable cover of $M$ by smooth charts. We need to show that for each smooth chart $\left( {V,\psi }\right)$ for $N$ , the set $\psi \left( {F\left( A\right)  \cap  V}\right)$ has measure zero in ${\mathbb{R}}^{n}$ . Note that this set is the union of countably many subsets of the form $\psi  \circ  F \circ  {\varphi }_{i}^{-1}\left( {{\varphi }_{i}\left( {A \cap  {U}_{i} \cap  {F}^{-1}\left( V\right) }\right) }\right)$ , each of which has measure zero by the result of Proposition 6.5.
\end{proof}


\section{The Whitney Approximation Theorems}

\begin{definition}
\label{def:6-41-extra-1}
\lean{delta_close}
\leanok
If $\delta  : M \rightarrow  \mathbb{R}$ is a positive continuous function, we say that two functions $F,\widetilde{F} : M \rightarrow  {\mathbb{R}}^{k}$ are $\delta$ -close if $\left| {F\left( x\right)  - \widetilde{F}\left( x\right) }\right|  < \; \delta \left( x\right)$ for all $x \in  M$ .
\end{definition}

\begin{corollary}
\label{cor:6-22}
\lean{Manifold.exists_positive_smooth_lt}
\leanok
\dcref{ch6:6.22}
If $M$ is a smooth manifold with or without boundary and $\delta  : M \rightarrow \; \mathbb{R}$ is a positive continuous function, there is a smooth function $e : M \rightarrow  \mathbb{R}$ such that $0 < e\left( x\right)  < \delta \left( x\right)$ for all $x \in  M$ .
\end{corollary}

\begin{proof}
\leanok
Use the Whitney approximation theorem to construct a smooth function $e : M \rightarrow  \mathbb{R}$ that satisfies $\left| {e\left( x\right)  - \frac{1}{2}\delta \left( x\right) }\right|  < \frac{1}{2}\delta \left( x\right)$ for all $x \in  M$ .
\end{proof}


\section{Tubular Neighborhoods}

\begin{definition}
\label{def:6-42-extra-1}
\lean{normal_bundle}
\leanok
For each $x \in  {\mathbb{R}}^{n}$ , the tangent space ${T}_{x}{\mathbb{R}}^{n}$ is canonically identified with ${\mathbb{R}}^{n}$ itself, and the tangent bundle $T{\mathbb{R}}^{n}$ is canonically diffeomorphic to ${\mathbb{R}}^{n} \times  {\mathbb{R}}^{n}$ . By virtue of this identification, each tangent space ${T}_{x}{\mathbb{R}}^{n}$ inherits a Euclidean dot product. Suppose $M \subseteq  {\mathbb{R}}^{n}$ is an embedded $m$ -dimensional submanifold. For each $x \in  M$ , we define the normal space to $M$ at $x$ to be the $\left( {n - m}\right)$ -dimensional subspace ${N}_{x}M \subseteq  {T}_{x}{\mathbb{R}}^{n}$ consisting of all vectors that are orthogonal to ${T}_{x}M$ with respect to the Euclidean dot product. The normal bundle of $M$ , denoted by ${NM}$ , is the subset of $T{\mathbb{R}}^{n} \approx  {\mathbb{R}}^{n} \times  {\mathbb{R}}^{n}$ consisting of vectors that are normal to $M$ :

\[
{NM} = \left\{  {\left( {x, v}\right)  \in  {\mathbb{R}}^{n} \times  {\mathbb{R}}^{n} : x \in  M, v \in  {N}_{x}M}\right\}  .
\]

There is a natural projection ${\pi }_{NM} : {NM} \rightarrow  M$ defined as the restriction to ${NM}$ of $\pi  : T{\mathbb{R}}^{n} \rightarrow  {\mathbb{R}}^{n}.$
\end{definition}

\begin{definition}
\label{def:6-42-extra-3}
\lean{TopologicalRetraction}
\leanok
A retraction of a topological space $X$ onto a subspace $M \subseteq  X$ is a continuous map $r : X \rightarrow  M$ such that ${\left. r\right| }_{M}$ is the identity map of $M$ .
\end{definition}


\section{Smooth Approximation of Maps Between Manifolds}

\begin{definition}
\label{def:6-43-extra-1}
\lean{ContMDiffMap.SmoothlyHomotopic}
\leanok
If $N$ and $M$ are two smooth manifolds with or without boundary, a homotopy $H : N \times  I \rightarrow  M$ is called a smooth homotopy if it is also a smooth map, in the sense that it extends to a smooth map on some neighborhood of $N \times  I$ in $N \times  \mathbb{R}$ . Two maps are said to be smoothly homotopic if there is a smooth homotopy between them.
\end{definition}

\begin{lemma}
\label{lem:6-28}
\lean{ContMDiffMap.SmoothlyHomotopic.equivalence}
\leanok
\dcref{ch6:6.28}
If $N$ and $M$ are smooth manifolds with or without boundary, smooth homotopy is an equivalence relation on the set of all smooth maps from $N$ to $M$ .
\end{lemma}

\begin{proof}
\leanok
Reflexivity and symmetry are proved just as for ordinary homotopy. To prove transitivity, suppose $F, G, K : N \rightarrow  M$ are smooth maps, and ${H}_{1},{H}_{2} : N \times  I \rightarrow  M$ are smooth homotopies from $F$ to $G$ and $G$ to $K$ , respectively. Let $\varphi  : \left\lbrack  {0,1}\right\rbrack   \rightarrow  \left\lbrack  {0,2}\right\rbrack$ be a smooth map such that $0 \leq  \varphi \left( t\right)  \leq  1$ for $t \in  \left\lbrack  {0,\frac{1}{2}}\right\rbrack  ,1 \leq  \varphi \left( t\right)  \leq  2$ for $t \in  \left\lbrack  {\frac{1}{2},1}\right\rbrack$ , $\varphi \left( 0\right)  = 0,\varphi \left( 1\right)  = 2$ , and $\varphi \left( t\right)  \equiv  1$ for $t$ in a neighborhood of $\frac{1}{2}$ (see Fig. 6.9).Define $H : N \times  I \rightarrow  M$ by

\[
H\left( {x, t}\right)  = \left\{  \begin{array}{ll} {H}_{1}\left( {x,\varphi \left( t\right) }\right) , & t \in  \left\lbrack  {0,\frac{1}{2}}\right\rbrack  , \\  {H}_{2}\left( {x,\varphi \left( t\right)  - 1}\right) , & t \in  \left\lbrack  {\frac{1}{2},1}\right\rbrack  . \end{array}\right.
\]

Then it is easy to check that $H$ is a smooth homotopy from $F$ to $K$ .
\end{proof}


\section{Transversality}

\begin{definition}
\label{def:6-44-extra-2}
\lean{IsSmoothFamily}
\leanok
Suppose $N, M$ , and $S$ are smooth manifolds, and for each $s \in  S$ we are given a map ${F}_{s} : N \rightarrow  M$ . The collection $\left\{  {{F}_{s} : s \in  S}\right\}$ is called a smooth family of maps if the map $F : M \times  S \rightarrow  N$ defined by $F\left( {x, s}\right)  = {F}_{s}\left( x\right)$ is smooth. You should think of such a family as a higher-dimensional analogue of a homotopy. The next proposition shows how such families are related to ordinary homotopies.
\end{definition}

\begin{proposition}
\label{prop:6-34}
\lean{chartedSpace_locPathConnectedSpace}
\leanok
\dcref{ch6:6.34}
If $\left\{  {{F}_{s} : s \in  S}\right\}$ is a smooth family of maps from $N$ to $M$ and $S$ is connected, then for any ${s}_{1},{s}_{2} \in  S$ , the maps ${F}_{{s}_{1}},{F}_{{s}_{2}} : N \rightarrow  M$ are homotopic.
\end{proposition}

\begin{proof}
\leanok
Because $S$ is connected, it is path-connected. If $\gamma  : \left\lbrack  {0,1}\right\rbrack   \rightarrow  S$ is any path from ${s}_{1}$ to ${s}_{2}$ , then $H\left( {x, s}\right)  = F\left( {x,\gamma \left( s\right) }\right)$ is a homotopy from ${F}_{{s}_{1}}$ to ${F}_{{s}_{2}}$ .
\end{proof}

\begin{definition}
\label{def:6-44-extra-3}
\lean{has_measure_zero_in_manifold}
\leanok
If $S$ is a smooth manifold and $B \subseteq  S$ is a subset whose complement has measure zero in $S$ , we say that $B$ contains almost every element of $S$ .
\end{definition}


\section{Problems}

\begin{exercise}
\label{prob:6-3}
\lean{Manifold.exists_smooth_zero_on_and_positive_off_lt_of_isClosed}
\leanok
\uses{cor:6-22}
6-3. Let $M$ be a smooth manifold, let $B \subseteq  M$ be a closed subset, and let $\delta  : M \rightarrow \; \mathbb{R}$ be a positive continuous function. Show that there is a smooth function $\widetilde{\delta } : M \rightarrow  \mathbb{R}$ that is zero on $B$ , positive on $M \smallsetminus  B$ , and satisfies $\widetilde{\delta }\left( x\right)  < \delta \left( x\right)$ everywhere. [Hint: consider $f/\left( {f + 1}\right)$ , where $f$ is a smooth nonnegative function that vanishes exactly on $B$ , and use Corollary 6.22.]
\end{exercise}

